%=============================================================
% Chapitre 7 — Interface applicative React 19
%=============================================================
\chapter{Interface applicative}
\label{ch:ui}

\section{Architecture technique frontend}

L'interface web Electio-Analytics est développée en \textbf{React 19} avec TypeScript.

\begin{table}[H]
\centering
\caption{Stack technique frontend}
\label{tab:stack}
\begin{tabular}{lll}
\toprule
\textbf{Composant} & \textbf{Bibliothèque} & \textbf{Rôle} \\
\midrule
Routing          & TanStack Router v1   & Navigation sans rechargement \\
Visualisation    & Recharts 2.x         & Graphiques interactifs \\
Animations       & Framer Motion 11     & Transitions fluides \\
Styles           & Tailwind CSS 4       & Design system utilitaire \\
État global      & TanStack Query       & Cache et synchronisation API \\
Données          & MongoDB + PyMongo    & Prédictions, statistiques, graphiques \\
\bottomrule
\end{tabular}
\end{table}

\section{Pages et composants principaux}

\subsection{Dashboard principal (\texttt{index.tsx})}

Le tableau de bord principal offre :
\begin{itemize}
  \item Un sélecteur de \textbf{modèle ML} (HistGradientBoosting, XGBoost, etc.)
  \item Des filtres géographiques \textbf{en cascade} : Région $\rightarrow$ Département $\rightarrow$ Canton
  \item Des cartes de KPIs (accuracy, F1-score, nombre de cantons prédits)
  \item Un graphique Recharts affichant les scores par candidat pour le territoire sélectionné
\end{itemize}

\subsection{Page Visualisation (\texttt{visualisation.tsx})}

La page intègre les graphiques Python (PNG générés par \texttt{generate\_visualisations.py}) :
\begin{itemize}
  \item Matrice de corrélations (heatmap)
  \item Comparaison des modèles (barres + courbes d'apprentissage)
  \item Importance des features (barres horizontales)
  \item Scatter emploi vs vote populiste
  \item Distribution des orientations par département
  \item Projections temporelles 2022--2026
\end{itemize}

\subsection{Architecture de routage}

\medskip\noindent\textit{TanStack Router --- routes principales :}
\begin{verbatim}
const router = createBrowserRouter([
  { path: "/",             component: DashboardPage },
  { path: "/visualisation", component: VisualisationPage },
  { path: "/resultats/:dept", component: ResultatsDeptPage },
  { path: "/comparaison",  component: ComparaisonModelesPage },
]);
\end{verbatim}

\section{Flux de données}

Les données transitent depuis \textbf{MongoDB} vers le dashboard React via une
couche API Python (FastAPI / Flask) qui expose les collections MongoDB en JSON :

\begin{center}
\begin{tabular}{ccccc}
\fbox{MongoDB \texttt{predictions}} & $\longrightarrow$ & \fbox{API Python} & $\longrightarrow$ & \fbox{React Query} \\[0.3cm]
\fbox{MongoDB \texttt{tendances}}   & $\longrightarrow$ & \fbox{API Python} &                   & $\downarrow$ \\[0.3cm]
\fbox{MongoDB \texttt{charts}}      & $\longrightarrow$ & \fbox{API Python} & $\longrightarrow$ & \fbox{Composants Recharts} \\
\end{tabular}
\end{center}

\medskip\noindent
Les collections MongoDB exposées par l'API backend :
\begin{itemize}
  \item \texttt{predictions} --- scores par candidat, département et modèle
  \item \texttt{tendances} --- projections temporelles 2022--2026
  \item \texttt{models\_meta} --- métriques et hyperparamètres des 5 modèles
  \item \texttt{charts} --- images de visualisation (PNG encodées en base64)
\end{itemize}

\section{Expérience utilisateur}

\begin{itemize}
  \item \textbf{Progressive disclosure} : les détails apparaissent au clic/survol
  \item \textbf{Couleurs sémantiques} : bleu = Centre, rouge = Extrême-Droite,
    rose = Gauche, orange = Droite, violet = Extrême-Gauche
  \item \textbf{Animations Framer Motion} : transitions 300ms ease-in-out
  \item \textbf{Responsive design} : adaptation mobile/tablette/desktop
\end{itemize}

\section{Génération des visualisations Python}

Les graphiques sont produits depuis les données de la collection \texttt{data\_final}
de MongoDB, puis les images PNG sont stockées dans la collection \texttt{charts}
(encodage base64), ce qui permet au dashboard de les récupérer sans accès
au système de fichiers local :

\medskip\noindent\textit{Extrait de \texttt{generate\_visualisations.py} :}
\begin{verbatim}
import matplotlib.pyplot as plt, seaborn as sns
import base64, io
from pymongo import MongoClient

db = MongoClient("mongodb://localhost:27017/")["electio_analytics"]
df = pd.DataFrame(list(db["data_final"].find({}, {"_id": 0})))

# Génération de la heatmap
fig, ax = plt.subplots(figsize=(10, 8))
sns.heatmap(corr_matrix, annot=True, fmt=".2f", ax=ax)

# Encodage en base64 et stockage dans MongoDB
buf = io.BytesIO()
plt.savefig(buf, format="png", dpi=150, bbox_inches="tight")
img_b64 = base64.b64encode(buf.getvalue()).decode()
db["charts"].update_one(
    {"name": "correlation_heatmap"},
    {"$set": {"data": img_b64, "format": "png"}},
    upsert=True
)
plt.close()
\end{verbatim}

\noindent Le même patron est appliqué aux six graphiques du projet :
heatmap, comparaison des modèles, importance des features, scatter emploi/vote,
distribution des orientations, projections temporelles.
